% abstract_v1.tex -- C9 paper lane (P9-01): abstract shared by the public and the anonymous build.
% Exponent 1/2 restricted to p(x)=x^{q-2}, odd m, d(a)=1+(a/rho)^4+(a/rho)^2, rho!=0 (finding C8-F02); the general-d
% repair statement keeps "every admissible d"; q=2^m and N=q^2 are stated.
We study a two-digit coordinate model of maps on the Galois ring $\GR(4,m)$, with $q=2^m$ and $N=q^2=|\GR(4,m)|$, in which the
ring addition carries from the low digit into the high digit. For maps
$F(a,b)=(p(a),\,h(a)\oplus d(a)b)$ with $p$ an APN permutation, $d$ a nowhere-vanishing
$\FF$-affine map whose linear part has a one-dimensional kernel, and $h$ arbitrary, we prove
direction-wise differential counts for $B=J\circ F$, a key-averaged two-block bound of order $1/N$, and a
collision-energy bound, all uniform in $h$. In contrast, a local integral property of the same
two-block map depends on $h$; for the inverse map and every admissible $d$ we give a linear
certificate on $h$ that excludes a family of local zero sums, and an explicit correction of Boolean
degree two of every $h$ that fails it; degree two is minimal for this local problem. For the averaged
alternating two-word arrangement we prove a one-word contraction of order $q^{-(1-\frac12\log_2 3)}$, uniform in $h$ and in the
admissible data, and deduce that a fixed number of steps brings the averaged two-word difference distribution to the scale $1/N$.
For the inverse map with $d(a)=1+(a/\rho)^4+(a/\rho)^2$, where $m$ is odd and $\rho\ne0$, the contraction improves to the order
$q^{-1/2}$, which is optimal since no admissible data give a faster order; for other $d$ our constant grows with the linearized
degree of $d$.
We then record what these statements do and do not imply for an independent-key composition:
the averaged multi-word bound coexists with a candidate-existence distinguisher with data linear
in the word size. No security claim for a concrete cipher is made.
